(1) 扭转运动方程 $I\ddot\theta+b\dot\theta+\kappa\theta=\tau_0\cos\omega t$。令 $\theta=\Re[\Theta e^{-i\omega t}]$，则 $(-I\omega^2-ib\omega+\kappa)\Theta=\tau_0$，所以复响应 $\Theta=\tau_0/(\kappa-I\omega^2-ib\omega)$，其分母包括惯性、耗散及弹性项。
(2) 稳态振幅为 $A=|\Theta|=\tau_0/\sqrt{(\kappa-I\omega^2)^2+b^2\omega^2}$。写 $\theta=A\cos(\omega t-\delta)$，滞后角可取 $\delta=\operatorname{atan2}(b\omega,\kappa-I\omega^2)$，在 $b,\omega>0$ 时介于零和 $\pi$，且共振附近接近 $\pi/2$。
(3) 令分母平方 $D(\omega)=(\kappa-I\omega^2)^2+b^2\omega^2$，对非零 $\omega$ 求极值：$D'=2\omega[-2I(\kappa-I\omega^2)+b^2]=0$。因此位移峰角频率 $\omega_{\rm pk}^2=\kappa/I-b^2/(2I^2)$；仅当 $\kappa/I>b^2/(2I^2)$ 才存在正频率振幅峰，否则最大值在静态 $\omega=0$。
(4) 阻尼瞬时热功率 $P_{\rm diss}=b\dot\theta^2$，时间平均 $\overline P_{\rm diss}=bA^2\omega^2/2$。驱动平均功率 $\overline{\tau(t)\dot\theta(t)}=\tau_0A\omega\sin\delta/2=bA^2\omega^2/2$，两者一致；$\omega\to0$ 时 $A\to\tau_0/\kappa$，$\omega\to\infty$ 时 $A\sim\tau_0/(I\omega^2)$。
